\section{The Speculative Speculative Decoding Framework}
\label{sec:ssd}

We introduce \textit{speculative speculative decoding}, a framework to reason about asynchronous variants of speculative decoding (Section~\ref{sec:ssd_alg}), and present results comparing the expected speed of SSD to baselines (Section~\ref{sec:ssd_theory}). \textbf{Definitions introduced here will be important in Section \ref{sec:method}.}

\subsection{Algorithm}
\label{sec:ssd_alg}

We present the SSD framework in Algorithm~\ref{alg:ssd}.
The speculator and verifier processes run in parallel on separate hardware.
While the verifier is verifying the drafted tokens from round $T$, the speculator begins speculating round $T+1$. 
It does so by predicting likely verification outcomes, and then preparing speculations for each of these outcomes in parallel, and storing these in a ``speculation cache'' (defined formally below).
Then, when it receives the actual verification outcome, it checks whether this was one of the outcomes in the cache that it had prepared for. If so, it immediately returns it.
Otherwise it defers to a fallback speculation strategy. 

For intuition on the SSD algorithm and why it is lossless, consider the similarities with speculative execution (\cite{specexec97}) for CPU optimization. 
Like speculative execution---which uses idle compute to pre-execute conditional paths in case they are later needed---SSD pre-speculates over many possible verification outcomes so that, when one occurs, the corresponding token sequence is immediately available.
That sequence is then verified exactly as in SD.
If the realized outcome was not pre-speculated, SSD falls back to a synchronous speculator, which is lossless since it reduces to ordinary SD. The following definitions will be important in understanding the SSD framework.

\begin{algorithm}[!ht]

\DontPrintSemicolon
\SetAlgoLined
\SetKwFunction{Main}{main}
\SetKwFunction{Verifier}{verifier}
\SetKwFunction{Speculator}{speculator}
\SetKwFunction{Prefill}{prefill}
\SetKwFunction{Verify}{verify}
\SetKwFunction{Speculate}{speculate}
\SetKwFunction{BuildCache}{build\_cache}
\SetKwFunction{PredictVerificationOutcomes}{predict\_verify\_outcomes}
\SetKwFunction{SpeculateOnOutcomes}{speculate\_for\_outcomes}
\SetKwFunction{Fallback}{fallback\_speculate}
\SetKwProg{Fn}{Function}{:}{}

\Fn{\Main{\text{prompt, target, primary\_draft, backup\_draft}}}{
    asynchronously launch \Speculator{prompt, primary\_draft, backup\_draft}\;
    $generated\_tokens \gets$ \Verifier{prompt, target}\;
    \Return $generated\_tokens$\;
}

\BlankLine

\Fn{\Verifier{prompt, target}}{
    $target$.\Prefill{prompt}\;
    \textbf{WAIT TO RECEIVE} $spec\_tokens$ from \Speculator\;
    $generated\_tokens \gets []$\;
    \While{True}{
        $verify\_outcome \gets target.\Verify{spec\_tokens}$\;
        $generated\_tokens$.append($verify\_outcome.tokens$)\;
        \textbf{SEND} $verify\_outcome$ to \Speculator\;
        \If{end\_token $\in verify\_outcome$}{
            \Return $generated\_tokens$\;
        }
        \textbf{WAIT TO RECEIVE} $spec\_tokens$ from \Speculator\;
    }
}

\BlankLine

\Fn{\Speculator{prompt, primary\_draft, backup\_draft}}{
    $primary\_draft$.\Prefill{prompt}\;
    $spec\_tokens \gets primary\_draft.\Speculate{prompt}$\;
    \While{True}{
        \textbf{SEND} $spec\_tokens$ to \Verifier for verification\;
        % $cache \gets$ \BuildCache{spec\_tokens} // see Sections~4.2--4.3\;

        \textcolor{blue}{$outcomes \gets$ \PredictVerificationOutcomes{spec\_tokens, primary\_draft} // Sec.~\ref{sec:cache_topology}}
        
        \textcolor{blue}{$cache \gets$ \SpeculateOnOutcomes{outcomes, primary\_draft} // Sec.~\ref{sec:cache_sampling}}
        
        
        \textbf{WAIT TO RECEIVE} $verify\_outcome$ from \Verifier\;
        \If{end\_token $\in verify\_outcome$}{
            \Return
        }
        \eIf{$verify\_outcome \in cache$}{
    $spec\_tokens \gets cache[verify\_outcome]$\;
}{
     \textcolor{blue}{
        $spec\_tokens \gets \Fallback$
        (\\ \quad \quad verify\_outcome, primary\_spec, backup\_spec\\
        )  // Sec.~\ref{sec:cache_misses}
    }\;
}

    }
}

\caption{The Speculative Speculative Decoding (SSD) Framework.}
\label{alg:ssd}
\end{algorithm}


\begin{definition}
    We define the \textbf{speculation cache} $\Spec^T$ for round $T$ as a dictionary mapping from a set of verification outcomes $\Ver^T$ to precomputed speculations for those outcomes.
    We denote a speculated token sequence $\spec^{T}$ corresponding to an outcome $\ver^T \in \Ver^T$ by $\Spec^T(\ver^{T})$ as a cache lookup operation.
\end{definition}

\begin{definition}
A \textbf{cache hit} is when the outcome $\ver^T$ of verifying $\spec^{T-1}$ is contained in the speculation cache $\Spec^T$.
A \textbf{cache miss} is when $\ver^T \notin \Spec^T$.
\end{definition}

\begin{definition}
    Let $\mathbf{p_{\mathrm{hit,p}}}$ denote the probability of a cache hit on an iteration conditional on the previous iteration having been speculated by the \textit{primary} speculator. Analogously, $\mathbf{p_{\mathrm{hit, b}}}$ represents the probability of a cache hit on an iteration conditional on speculating with the \textit{backup} model. Finally, let $\mathbf{p_{\mathrm{hit}}}$ denote the \textit{unconditional} probability of a cache hit in an iteration. 
\end{definition}



\subsection{Theoretical Results}
\label{sec:ssd_theory}
We analyze the performance of \Sys relative to regular autoregressive decoding (Theorem~\ref{thm:speedup}) and ordinary speculative decoding (Corollary~\ref{cor:speedup}).
Proofs deferred to Appendix~\ref{app:theory}.

\begin{theorem}
\label{thm:speedup}
Let the primary and backup speculators take time $T_p$ and $T_b$ relative to the verifier. Let $E_{\mathrm{hit}}$ and $E_{\mathrm{miss}}$ denote the expected number of generated tokens from the primary and backup speculators, respectively. The expected speedup of Algorithm~\ref{alg:ssd} relative to autoregressive decoding is then: 
\begin{eqnarray*}
speedup_{\mathrm{SSD}} &=& \frac{p_{\mathrm{hit}} \cdot E_{\mathrm{hit}} + (1-p_{\mathrm{hit}}) \cdot E_{\mathrm{miss}}}{p_{\mathrm{hit}} \cdot \max(1, T_p) + (1-p_{\mathrm{hit}}) \cdot (1 + T_b)}. \\
\end{eqnarray*}

% Further, define $p_{\mathrm{hit,p}}$ and $p_{\mathrm{hit,b}}$ as the cache hit rates conditioning on the previous speculation being from the primary or backup speculator, respectively. Then, assuming $|p_{\mathrm{hit,p}}-p_{\mathrm{hit,b}}|<1$, the global cache hit rate $p_{\mathrm{hit}}$ is given by: 
% \begin{eqnarray*}
% p_{\mathrm{hit}} &=& \frac{p_{\mathrm{hit,b}}}{1 + p_{\mathrm{hit,b}} - p_{\mathrm{hit,p}} }.
% \end{eqnarray*}
\end{theorem}
The numerator corresponds to the expected number of tokens generated in each iteration of the algorithm, and the denominator corresponds to the expected latency of each iteration relative to autoregressive decoding. This implies two key corollaries. 

\begin{corollary} (\textbf{Strictly Faster Than SD})
\label{cor:stricly_faster}
    Suppose we run SD with a given speculator $\mathcal{M}$.
    Running SSD with primary and backup both set to $\mathcal{M}$ does no worse than SD (strictly better if $p_{\mathrm{hit}} > 0$, $T_{\mathrm{SD}} > 0$). 
\end{corollary}

% Thus, in the worst case when we have only cache misses, and fall back to using the same speculator we would in SD, SSD reduces to SD and does no worse. 

\begin{corollary} (\textbf{Speedup Sandwich})
\label{cor:speedup}
Suppose we choose a primary speculator for which drafting completes before verification ($T_p < 1$), and a fast backup speculator ($T_b = 0$).
Then if $T_{\mathrm{SD}}, E_{\mathrm{SD}}$ represent the latency and expected number of generated tokens from a draft model in SD, then the SSD speedup over SD can be bounded by:
\begin{eqnarray*}
\Big(1 + T_{\mathrm{SD}}\Big) \cdot \frac{E_{\mathrm{hit}}}{E_{\mathrm{SD}}} \cdot p_{\mathrm{hit}} \leq \frac{speedup_{\mathrm{SSD}}}{speedup_{\mathrm{SD}}} \leq \Big(1 + T_{\mathrm{SD}}\Big) \cdot \frac{E_{\mathrm{hit}}}{E_{\mathrm{SD}}}.
\end{eqnarray*}
\end{corollary}
This equation reveals that the maximum speedup attainable by SSD is proportional to the latency reduction $(1+T_{\mathrm{SD}})$ from hiding draft latency, and the increase in expected number of generated tokens ($E_{\mathrm{hit}}/E_{\mathrm{SD}}$) from increased drafting time.
However, a low cache hit rate $p_{\mathrm{hit}}$ reduces the effectiveness of this algorithm, as shown by the lower bound. 


% Special cases:
% 1. p_hit=0, E_miss = E_SD
% ====> This is equivalent to sequenial SD.
% 2. p_hit > 0, E_hit > E_SD (:=E_miss), max(1,T_p) < 1 + T_SD.
% ====> SSD >= SD.
% 3. ASSUMPTION #1: max(1,T_p) = 1, T_b \approx 0
% ====> simplifies to p_hit * E_hit + (1-p_hit) * E_miss
% 4. Sandwich relative to SD. uses assumption 1.
% ====> (p_hit * E_hit + (1-p_hit) * E_miss) / (E_SD / (1+T_SD))

% E_miss = E_SD


% Assumptions #1:
% Setting your backup to SD spec --> greatly better
% p_hit >= 0, E_hit >= E_SD (:=E_miss), max(1,T_p) <= 1 + T_SD.
% ==> speedup SSD >= SD

% Assumptions #2
% -- max(1,T_p) = 1, T_b \approx 0 (super fast backup, fast enough primary).
% -- p_hit = 1 gives upper bound, p_hit < 1 gives lower



% T_p < 1, T_b = 0, p_hit = 1 ===> upper bound




% p_hit * E_hit + (1-p_hit) * E_SD / (p_hit * max(1,T_p) + (1-p)*(1+T_b))
% >> numerator > E_SD
% >> denominator < 1+T_SD

% E_SD / 


% vs
% E_SD / (1 + T_SD)



% In practice, we choose speculators and lookahead length, as well as the backup speculator\footnote{There exist extremely fast backup speculators, such as speculating a constant sequence, random tokens, or using n-gram based (non-neural) language models~\citep{infinigram,suffixtree}}, so that the following assumption is satisfied:

% \textbf{Assumption 1.}
% We assume that speculation always completes before verification so that $T_p < 1$ and $T_b \approx 0$, thus the latency of drafting is hidden.
% Each iteration of the algorithm therefore takes the same amount of time as a single forward pass of the target model, so the speedup equation reduces to:

% % \TK{where do we use this assumption in the paper? directly below?}
% \begin{eqnarray*}
% speedup &=& p_{\mathrm{hit}} \cdot E_{\mathrm{hit}} + \big(1-p_{\mathrm{hit}}\big) \cdot E_{\mathrm{miss}}.
% \end{eqnarray*}

